\documentclass[10pt,a4paper]{amsart}
\usepackage[margin=19mm]{geometry}
\usepackage{amsmath,amssymb,mathtools}
\usepackage[scr=rsfs,cal=euler]{mathalfa}
\usepackage{xfrac}
\usepackage[unicode,hidelinks,bookmarks=false]{hyperref}
\newcommand{\ie}{i.e.\ }
\newcommand{\cf}{cf.\ }
\newcommand{\resp}{respectively\ }
\newcommand{\Subsection}{\subsection}
\providecommand{\nobreakdash}{\nobreak\hbox{-}}
\setlength{\emergencystretch}{2em}
\multlinegap0pt
\usepackage{amssymb}
\usepackage{mathtools}
\usepackage[shortlabels]{enumitem}
\newtheorem{proposition}{Proposition}
\newcommand{\C}{\mathbf C}
\newcommand{\PP}{\mathbf P}
\newcommand{\Pic}{\operatorname{Pic}}
\newcommand{\Z}{\mathbf Z}
\newcommand{\cO}{\mathcal O}
\newcommand{\divisor}{\operatorname{div}}
\title{Degree of irrationality of abelian surfaces}
\author{Songsong Huang, Zhi Jiang, Xin Yang}
\date{Source: arXiv:2609.02169v1 (CC BY 4.0)}
\begin{document}
\maketitle

\noindent\emph{Source and licence.} Degree of irrationality of abelian surfaces. Version arXiv:2609.02169v1 (CC BY 4.0), distributed by arXiv under the Creative Commons Attribution 4.0 International licence, http://creativecommons.org/licenses/by/4.0/, as recorded in the arXiv licence metadata for this exact version and shown on the abstract page https://arxiv.org/abs/2609.02169v1. Full licence notice: \texttt{licenses/arxiv-2609.02169-CC-BY-4.0.txt}.\par\smallskip

\noindent\textit{Editorial scope.} Counted unit: the article's proposition prop:good-theta (source0.tex lines 1425--1432). The article's complete original proof of it is delivered with it (source0.tex lines 1434--1554, one \texttt{proof} environment of 121 lines). No mathematics is written, repaired or re-derived: the statement and the proof are the article's own lines, in the article's own notation, and no retained line is substituted or re-flowed.  Retained uncounted as the source-local context the proof needs: lines 272--275.  Nothing was omitted from any retained span, so the package carries no omission record.  No retained line contains a citation, so the wrapper carries no bibliography and no bibliography entry.  No retained line contains a figure environment, an image-inclusion command or drawing code, so no image is delivered and the package declares no asset dependency.  Editorial formatting only, in addition: an AMS \texttt{amsart} preamble that restates the article's own declarations and macros used by the retained lines, namely the environments \texttt{proposition} on separate counters, together with the standard packages those lines need.  The retained environments print in this document as Proposition 1 in order of appearance, whereas the article prints the same items with the numbering of its own full document.  The complete native source of the article as distributed in the arXiv e-print for this exact version is bundled unchanged as \texttt{sources/209180/source0.tex}.
\begin{equation}
  \cO_A(D+x)\simeq T_{-x}^*\cO_A(D).
\label{eq:translate-convention}
\end{equation}

\begin{proposition}\label{prop:good-theta}
Let $(B,\Theta)$ be an indecomposable principally polarized complex
abelian surface.  Let $C\subset B$ be a smooth symmetric theta divisor,
and let $0\neq t\in B[3]$ satisfy $t\notin[2]_B(C)$.  Then there is a
dominant rational map $f_t:B\dashrightarrow\PP^2$
of degree three which is equivariant for the translation
$T_t:B\to B$ and a faithful projective action of order three on $\PP^2$.
\end{proposition}

\begin{proof}
Choose $p\in C\cap(C+t)$ and put
\[
  q=p-t,\qquad r=p+t=q-t,\qquad d=p+q.
\]
Then $p,q\in C$.  Moreover, the element $d$ satisfies
\begin{equation}
  d\notin\{0,t,-t\},
\label{eq:d-not}
\end{equation}
because these three equalities would respectively give
$t=2p$, $t=2(-p)$, or $t=2q$, contradicting
$t\notin[2]_B(C)$.

With indices in $\Z/3\Z$, define
\[
  B_i=C+it,\qquad A_i=C+d+it.
\]
Since $C$ represents a principal polarization,
$\operatorname{Stab}(C)\subseteq K(\cO_B(C))=\{0_B\}$.
It follows that the three $B_i$ are pairwise distinct, as are the three
$A_i$.  If $A_i=B_j$, then
$d+(i-j)t\in\operatorname{Stab}(C)$, and hence
$d\in\langle t\rangle=\{0,t,-t\}$, contradicting
\eqref{eq:d-not}.  Thus all six curves are pairwise distinct.
Let $b_i$ and $a_i$ be canonical sections cutting out $B_i$ and $A_i$.
Put $L_i=\cO_B(B_i)$ and
$Q=\phi_{\cO_B(C)}(d)\in\Pic^0(B)$.  By
\eqref{eq:translate-convention} and the theorem of the square,
$\cO_B(A_i)\simeq L_i\otimes Q^{-1}$.
Consequently, the products
\[
  s_i:=a_i b_{i+1}b_{i+2}\qquad(i\in\Z/3\Z)
\]
are sections of the same line bundle
\[
  M=(L_0\otimes L_1\otimes L_2)\otimes Q^{-1}.
\]
Numerically,
\[
  c_1(M)=3[\Theta],\qquad M^2=18.
\]
The sections are linearly independent: restricting a relation to $B_i$
leaves only the restriction of $s_i$, which is not identically zero.

We use the following intersection rule.  If $x,y\in C$ and
$C\neq C+x+y$, then, scheme-theoretically,
\begin{equation}
 C\cap(C+x+y)=
 \begin{cases}
   \{x,y\},&x\neq y,\\
   2[x],&x=y.
 \end{cases}
\label{eq:theta-intersection}
\end{equation}
Indeed, $C=-C$, so the displayed points lie in the intersection, whose
total length is $C^2=2$.  In the double case the involution
$z\mapsto2x-z$ identifies the two tangent lines at $x$.
Applying \eqref{eq:theta-intersection} gives
\begin{equation}
\begin{array}{c|ccc}
 &B_0&B_1&B_2\\ \hline
A_0&\{p,q\}&2[p]&2[q]\\
A_1&2[p]&\{p,r\}&2[r]\\
A_2&2[q]&2[r]&\{q,r\}
\end{array}
\label{eq:A-B-table}
\end{equation}
and
\begin{equation}
\begin{aligned}
 B_0\cap B_1&=\{p,-q\},\\
 B_0\cap B_2&=\{q,-p\},\\
 B_1\cap B_2&=\{r,-r\}.
\end{aligned}
\label{eq:B-B-table}
\end{equation}
The six points
\begin{equation}
  p,q,r,-p,-q,-r
\label{eq:six-theta-points}
\end{equation}
are pairwise distinct.  Indeed, the positive points form the
$\langle t\rangle$-orbit of $p$; if this orbit met its negative, then
$2p\in\langle t\rangle$, and one of $p,-p,p-t\in C$ would have double
$t$, again contradicting goodness.

Equations \eqref{eq:A-B-table}--\eqref{eq:B-B-table} show that
\eqref{eq:six-theta-points} is the complete base support of the linear system.
A point on at least two $B_i$ is a base point.  A point on exactly one
$B_i$ is a base point precisely when it also lies on the corresponding
$A_i$.  A point on none of the $B_i$ would lie on all three $A_i$; its
three-point $t$-orbit would then lie in $A_0\cap A_1$, contradicting
$A_0\cdot A_1=2$.

At each of $p,q,r$, the exact base ideal is $\mathfrak m^2$.  For example,
at $p$, take regular parameters $x,y$ cutting out $B_0,B_1$.  Up to
units,
$(s_0,s_1,s_2)\equiv(y^2,x^2,xy)\pmod{\mathfrak m_p^3}$.
Thus the base ideal $I_p$ satisfies
$I_p+\mathfrak m_p^3=\mathfrak m_p^2$ and
$I_p\subseteq\mathfrak m_p^2$; Nakayama's lemma applied to
$\mathfrak m_p^2/I_p$ gives $I_p=\mathfrak m_p^2$.
At each of $-p,-q,-r$, exactly two $B_i$ meet transversely and no $A_i$
passes through the point, so the exact base ideal is $\mathfrak m$.

Let $\mu_t:X_t\to B$ be the blow-up of the six points, and write $E_z$
for the exceptional curve over $z$.  The moving divisor is
$H_t=\mu_t^*M-2(E_p+E_q+E_r)
-(E_{-p}+E_{-q}+E_{-r})$.  It is base-point-free and
$H_t^2=18-3\cdot2^2-3\cdot1^2=3$.
The induced morphism $X_t\to\PP^2$ is dominant, since a globally generated
divisor defining a map to a curve has square zero, and its degree is
$H_t^2=3$.

Finally, set $D_i=A_i+B_{i+1}+B_{i+2}=\divisor(s_i)$.
Pullback by $T_t$ gives $T_t^*D_i=D_{i-1}$.  Hence translation by $t$
cyclically permutes the three coordinate lines of the linear system, up to nonzero
scalars.  The resulting weighted three-cycle has exact order three in
$\operatorname{PGL}_3(\C)$, so the target action is faithful.
\end{proof}

\end{document}
